% --- Archon physics formalization source begin ---
% archon:physics
% archon:covers PhyXMiniProblems/problem_phyx_mini_0201.lean
% archon:source-report reports/phyx_mini/problem_phyx_mini_0201.source.json

\chapter{Physics problem phyx\_mini\_0201}
\label{ch:PhyXMiniProblems_problem_phyx_mini_0201}

\paragraph{Problem source.}
\#\# Physical scenario

Echolocation is a form of sensory perception used by animals such as bats, dolphins, and toothed whales. The animal emits a pulse of sound (a longitudinal wave) which, after reflection from objects, returns and is detected by the animal. Echolocation waves can have frequencies of about 100\{,\}000\textbackslash{},\textbackslash{}mathrm\{Hz\}.An obstacle is 100\textbackslash{},\textbackslash{}mathrm\{m\} from the animal.

\#\# Question

How long after the animal emits a wave is its reflection detected?

\#\# Answer choices

- A: A: 0.20s

- B: B: 0.18s

- C: C: 0.16s

- D: D: 0.14s

\#\# Figure caption (auxiliary; use the image as primary evidence)

The image shows a beluga whale underwater, swimming towards the camera. It is centered in the frame, and its mouth is slightly open. Above the whale, there is a large air bubble or a splash effect, creating a ring on the water surface. The background is a deep blue, typical of underwater scenes, with light reflecting from the surface. There is no text present in the image.

\#\# Dataset metadata

Wave Properties; Physical Model Grounding Reasoning, Predictive Reasoning

\paragraph{Recorded answer/context.}
D: D: 0.14s

\paragraph{Figure/image path.}
/root/proposal\_for\_physic/science-mango/phyx\_mini\_run/phyx\_data/test\_image/201.png

\paragraph{Formalization target.}
create a compiling Lean file with sorry bodies at `PhyXMiniProblems/problem\_phyx\_mini\_0201.lean`. The Lean declarations must preserve the physical quantities, dimensions or dimensional roles, figure labels, governing-law hypotheses, and final relation expressed by this problem.
Use Mathlib/Physlib names found through LeanExplore where available. If a physics API is missing, introduce faithful local abstractions rather than scalar placeholder aliases.

\begin{theorem}[Physics formalization target]
\label{thm:physics:phyx_mini_0201:target}
\lean{PhyXMiniProblems.ProblemPhyXMini0201.reflectionDetectionDelay_matches_recordedAnswerD}
\uses{def:physics:phyx-mini-0201:phyxminiproblems-problemphyxmini0201-echolocationsetup, def:physics:phyx-mini-0201:phyxminiproblems-problemphyxmini0201-matchesproblemandfigure, def:physics:phyx-mini-0201:phyxminiproblems-problemphyxmini0201-usestextbookwatersoundspeed, def:physics:phyx-mini-0201:phyxminiproblems-problemphyxmini0201-hasphysicalecholocationparameters, def:physics:phyx-mini-0201:phyxminiproblems-problemphyxmini0201-satisfiesroundtripsoundpropagation, def:physics:phyx-mini-0201:phyxminiproblems-problemphyxmini0201-matchesanswerchoice, def:physics:phyx-mini-0201:phyxminiproblems-problemphyxmini0201-isclosestanswerchoice, def:physics:phyx-mini-0201:phyxminiproblems-problemphyxmini0201-recordedanswerchoice}
The assigned autoformalize agent should translate this physics problem into Lean declarations in the covered file, with theorem and lemma proof bodies written as `by sorry`.
\end{theorem}
\begin{proof}
This is an autoformalization task, not a proof task. Produce statements that can later be proved by the physics prover without weakening the physical model.
\end{proof}

% --- Archon declaration topology begin ---
\section{Declaration topology}
\begin{definition}[Acoustic Length]
  \label{def:physics:phyx-mini-0201:phyxminiproblems-problemphyxmini0201-acousticlength}
  \lean{PhyXMiniProblems.ProblemPhyXMini0201.AcousticLength}
  A nonnegative physical acoustic-path length
\end{definition}
\begin{proof}
  This is the declared model definition of Acoustic Length.
\end{proof}
\begin{definition}[Acoustic Duration]
  \label{def:physics:phyx-mini-0201:phyxminiproblems-problemphyxmini0201-acousticduration}
  \lean{PhyXMiniProblems.ProblemPhyXMini0201.AcousticDuration}
  A nonnegative physical propagation duration
\end{definition}
\begin{proof}
  This is the declared model definition of Acoustic Duration.
\end{proof}
\begin{definition}[Acoustic Frequency]
  \label{def:physics:phyx-mini-0201:phyxminiproblems-problemphyxmini0201-acousticfrequency}
  \lean{PhyXMiniProblems.ProblemPhyXMini0201.AcousticFrequency}
  A nonnegative physical acoustic frequency, carrying inverse-time dimension
\end{definition}
\begin{proof}
  This is the declared model definition of Acoustic Frequency.
\end{proof}
\begin{definition}[length In Meters]
  \label{def:physics:phyx-mini-0201:phyxminiproblems-problemphyxmini0201-lengthinmeters}
  \lean{PhyXMiniProblems.ProblemPhyXMini0201.lengthInMeters}
  \uses{def:physics:phyx-mini-0201:phyxminiproblems-problemphyxmini0201-acousticlength}
  Read a physical acoustic-path length as a real number of metres
\end{definition}
\begin{proof}
  This is the declared model definition of length In Meters.
\end{proof}
\begin{definition}[duration In Seconds]
  \label{def:physics:phyx-mini-0201:phyxminiproblems-problemphyxmini0201-durationinseconds}
  \lean{PhyXMiniProblems.ProblemPhyXMini0201.durationInSeconds}
  \uses{def:physics:phyx-mini-0201:phyxminiproblems-problemphyxmini0201-acousticduration}
  Read a physical propagation duration as a real number of seconds
\end{definition}
\begin{proof}
  This is the declared model definition of duration In Seconds.
\end{proof}
\begin{definition}[frequency In Hertz]
  \label{def:physics:phyx-mini-0201:phyxminiproblems-problemphyxmini0201-frequencyinhertz}
  \lean{PhyXMiniProblems.ProblemPhyXMini0201.frequencyInHertz}
  \uses{def:physics:phyx-mini-0201:phyxminiproblems-problemphyxmini0201-acousticfrequency}
  Read a physical acoustic frequency as a real number of hertz
\end{definition}
\begin{proof}
  This is the declared model definition of frequency In Hertz.
\end{proof}
\begin{definition}[speed In Meters Per Second]
  \label{def:physics:phyx-mini-0201:phyxminiproblems-problemphyxmini0201-speedinmeterspersecond}
  \lean{PhyXMiniProblems.ProblemPhyXMini0201.speedInMetersPerSecond}
  Read Physlib's nonnegative dimensionful speed in metres per second
\end{definition}
\begin{proof}
  This is the declared model definition of speed In Meters Per Second.
\end{proof}
\begin{definition}[Echolocating Animal]
  \label{def:physics:phyx-mini-0201:phyxminiproblems-problemphyxmini0201-echolocatinganimal}
  \lean{PhyXMiniProblems.ProblemPhyXMini0201.EcholocatingAnimal}
  The animal species identified by the primary image
\end{definition}
\begin{proof}
  This is the declared model definition of Echolocating Animal.
\end{proof}
\begin{definition}[Acoustic Medium]
  \label{def:physics:phyx-mini-0201:phyxminiproblems-problemphyxmini0201-acousticmedium}
  \lean{PhyXMiniProblems.ProblemPhyXMini0201.AcousticMedium}
  The propagation medium visible in the primary image
\end{definition}
\begin{proof}
  This is the declared model definition of Acoustic Medium.
\end{proof}
\begin{definition}[Acoustic Wave Mode]
  \label{def:physics:phyx-mini-0201:phyxminiproblems-problemphyxmini0201-acousticwavemode}
  \lean{PhyXMiniProblems.ProblemPhyXMini0201.AcousticWaveMode}
  Sound is modeled as the longitudinal wave stated in the problem
\end{definition}
\begin{proof}
  This is the declared model definition of Acoustic Wave Mode.
\end{proof}
\begin{definition}[Echo Path Leg]
  \label{def:physics:phyx-mini-0201:phyxminiproblems-problemphyxmini0201-echopathleg}
  \lean{PhyXMiniProblems.ProblemPhyXMini0201.EchoPathLeg}
  The two legs of the pulse path before detection of the echo
\end{definition}
\begin{proof}
  This is the declared model definition of Echo Path Leg.
\end{proof}
\begin{definition}[Echolocation Setup]
  \label{def:physics:phyx-mini-0201:phyxminiproblems-problemphyxmini0201-echolocationsetup}
  \lean{PhyXMiniProblems.ProblemPhyXMini0201.EcholocationSetup}
  \uses{def:physics:phyx-mini-0201:phyxminiproblems-problemphyxmini0201-acousticlength, def:physics:phyx-mini-0201:phyxminiproblems-problemphyxmini0201-acousticduration, def:physics:phyx-mini-0201:phyxminiproblems-problemphyxmini0201-acousticfrequency, def:physics:phyx-mini-0201:phyxminiproblems-problemphyxmini0201-echolocatinganimal, def:physics:phyx-mini-0201:phyxminiproblems-problemphyxmini0201-acousticmedium, def:physics:phyx-mini-0201:phyxminiproblems-problemphyxmini0201-acousticwavemode, def:physics:phyx-mini-0201:phyxminiproblems-problemphyxmini0201-echopathleg}
  The physical quantities and labels of the echolocation experiment. reflectionDetectionDelay is an unknown physical duration. In particular, this structure does not assign it any displayed answer value
\end{definition}
\begin{proof}
  This is the declared model definition of Echolocation Setup.
\end{proof}
\begin{definition}[Matches Problem And Figure]
  \label{def:physics:phyx-mini-0201:phyxminiproblems-problemphyxmini0201-matchesproblemandfigure}
  \lean{PhyXMiniProblems.ProblemPhyXMini0201.MatchesProblemAndFigure}
  \uses{def:physics:phyx-mini-0201:phyxminiproblems-problemphyxmini0201-lengthinmeters, def:physics:phyx-mini-0201:phyxminiproblems-problemphyxmini0201-frequencyinhertz, def:physics:phyx-mini-0201:phyxminiproblems-problemphyxmini0201-echolocationsetup}
  Problem-text and primary-image data: a beluga in water emits a longitudinal 100000 Hz pulse toward an obstacle 100 m away. Reflection makes both path legs have the stated one-way length. No detection-time answer occurs here
\end{definition}
\begin{proof}
  This is the declared model definition of Matches Problem And Figure.
\end{proof}
\begin{definition}[Uses Textbook Water Sound Speed]
  \label{def:physics:phyx-mini-0201:phyxminiproblems-problemphyxmini0201-usestextbookwatersoundspeed}
  \lean{PhyXMiniProblems.ProblemPhyXMini0201.UsesTextbookWaterSoundSpeed}
  \uses{def:physics:phyx-mini-0201:phyxminiproblems-problemphyxmini0201-speedinmeterspersecond, def:physics:phyx-mini-0201:phyxminiproblems-problemphyxmini0201-echolocationsetup}
  The textbook water-sound-speed calibration implicit in the supplied answer table. This is auxiliary physical data, not the requested detection delay
\end{definition}
\begin{proof}
  This is the declared model definition of Uses Textbook Water Sound Speed.
\end{proof}
\begin{definition}[Has Physical Echolocation Parameters]
  \label{def:physics:phyx-mini-0201:phyxminiproblems-problemphyxmini0201-hasphysicalecholocationparameters}
  \lean{PhyXMiniProblems.ProblemPhyXMini0201.HasPhysicalEcholocationParameters}
  \uses{def:physics:phyx-mini-0201:phyxminiproblems-problemphyxmini0201-lengthinmeters, def:physics:phyx-mini-0201:phyxminiproblems-problemphyxmini0201-frequencyinhertz, def:physics:phyx-mini-0201:phyxminiproblems-problemphyxmini0201-speedinmeterspersecond, def:physics:phyx-mini-0201:phyxminiproblems-problemphyxmini0201-echolocationsetup}
  Strict positivity conditions for the emitted wave and propagation data
\end{definition}
\begin{proof}
  This is the declared model definition of Has Physical Echolocation Parameters.
\end{proof}
\begin{definition}[Satisfies Round Trip Sound Propagation]
  \label{def:physics:phyx-mini-0201:phyxminiproblems-problemphyxmini0201-satisfiesroundtripsoundpropagation}
  \lean{PhyXMiniProblems.ProblemPhyXMini0201.SatisfiesRoundTripSoundPropagation}
  \uses{def:physics:phyx-mini-0201:phyxminiproblems-problemphyxmini0201-lengthinmeters, def:physics:phyx-mini-0201:phyxminiproblems-problemphyxmini0201-durationinseconds, def:physics:phyx-mini-0201:phyxminiproblems-problemphyxmini0201-speedinmeterspersecond, def:physics:phyx-mini-0201:phyxminiproblems-problemphyxmini0201-echolocationsetup}
  Governing laws for the echo. Constant-speed travel holds on the outbound and reflected-return legs, and the detection delay is the sum of their travel times. These laws contain no numerical value for the requested delay
\end{definition}
\begin{proof}
  This is the declared model definition of Satisfies Round Trip Sound Propagation.
\end{proof}
\begin{definition}[Answer Choice]
  \label{def:physics:phyx-mini-0201:phyxminiproblems-problemphyxmini0201-answerchoice}
  \lean{PhyXMiniProblems.ProblemPhyXMini0201.AnswerChoice}
  Labels of the four answers displayed with the problem
\end{definition}
\begin{proof}
  This is the declared model definition of Answer Choice.
\end{proof}
\begin{definition}[seconds]
  \label{def:physics:phyx-mini-0201:phyxminiproblems-problemphyxmini0201-answerchoice-seconds}
  \lean{PhyXMiniProblems.ProblemPhyXMini0201.AnswerChoice.seconds}
  \uses{def:physics:phyx-mini-0201:phyxminiproblems-problemphyxmini0201-answerchoice}
  Seconds printed beside each displayed answer label
\end{definition}
\begin{proof}
  This is the declared model definition of seconds.
\end{proof}
\begin{definition}[Matches Answer Choice]
  \label{def:physics:phyx-mini-0201:phyxminiproblems-problemphyxmini0201-matchesanswerchoice}
  \lean{PhyXMiniProblems.ProblemPhyXMini0201.MatchesAnswerChoice}
  \uses{def:physics:phyx-mini-0201:phyxminiproblems-problemphyxmini0201-acousticduration, def:physics:phyx-mini-0201:phyxminiproblems-problemphyxmini0201-durationinseconds, def:physics:phyx-mini-0201:phyxminiproblems-problemphyxmini0201-answerchoice}
  A physical duration agrees with a displayed hundredth-second answer. The half-centisecond tolerance models rounding to the printed precision
\end{definition}
\begin{proof}
  This is the declared model definition of Matches Answer Choice.
\end{proof}
\begin{definition}[Is Closest Answer Choice]
  \label{def:physics:phyx-mini-0201:phyxminiproblems-problemphyxmini0201-isclosestanswerchoice}
  \lean{PhyXMiniProblems.ProblemPhyXMini0201.IsClosestAnswerChoice}
  \uses{def:physics:phyx-mini-0201:phyxminiproblems-problemphyxmini0201-acousticduration, def:physics:phyx-mini-0201:phyxminiproblems-problemphyxmini0201-durationinseconds, def:physics:phyx-mini-0201:phyxminiproblems-problemphyxmini0201-answerchoice}
  A displayed choice is uniquely closest to the physical duration
\end{definition}
\begin{proof}
  This is the declared model definition of Is Closest Answer Choice.
\end{proof}
\begin{definition}[recorded Answer Choice]
  \label{def:physics:phyx-mini-0201:phyxminiproblems-problemphyxmini0201-recordedanswerchoice}
  \lean{PhyXMiniProblems.ProblemPhyXMini0201.recordedAnswerChoice}
  \uses{def:physics:phyx-mini-0201:phyxminiproblems-problemphyxmini0201-answerchoice}
  The answer label recorded in the supplied dataset metadata
\end{definition}
\begin{proof}
  This is the declared model definition of recorded Answer Choice.
\end{proof}
\begin{lemma}[reflection Detection Delay exact]
  \label{lem:physics:phyx-mini-0201:phyxminiproblems-problemphyxmini0201-reflectiondetectiondelay-exact}
  \lean{PhyXMiniProblems.ProblemPhyXMini0201.reflectionDetectionDelay_exact}
  \uses{def:physics:phyx-mini-0201:phyxminiproblems-problemphyxmini0201-durationinseconds, def:physics:phyx-mini-0201:phyxminiproblems-problemphyxmini0201-echolocationsetup, def:physics:phyx-mini-0201:phyxminiproblems-problemphyxmini0201-matchesproblemandfigure, def:physics:phyx-mini-0201:phyxminiproblems-problemphyxmini0201-usestextbookwatersoundspeed, def:physics:phyx-mini-0201:phyxminiproblems-problemphyxmini0201-hasphysicalecholocationparameters, def:physics:phyx-mini-0201:phyxminiproblems-problemphyxmini0201-satisfiesroundtripsoundpropagation}
  Before rounding to the precision of the choices, the ideal round-trip model gives an exact detection delay of 1/7 s
\end{lemma}
\begin{proof}
  The conclusion follows from the governing relations and figure readouts named in the statement of reflection Detection Delay exact.
\end{proof}
% --- Archon declaration topology end ---
% --- Archon physics formalization source end ---
